\section{C2/C3 Public-Sample Surface Screen}
\label{app:c2-public-surface}

C2 and C3 reduce final compaction noise by replacing coordinate-wise compaction with a block dictionary.  The security cost is a much larger public q-ary LPN surface.  This appendix records the screening model used by the validation repository.  It is not a certified cryptanalytic estimator; it is a conservative design screen intended to expose surfaces that must be analyzed before any concrete security claim.

\paragraph{Dictionary row count.}
For block widths $\ell_j$, the C2 dictionary publishes
\[
  E = \sum_j (q^{\ell_j}-1)
\]
CLPN ciphertexts.  Each ciphertext contains $m_c$ noisy rows under the same compaction secret $r\in\mathbb F_q^{n_c}$, so the direct dictionary surface contains $E m_c$ public rows.  Seeded C3 changes only the representation of the public matrices: it stores a seed for each matrix and the corresponding $b$ vector.  Since the adversary can regenerate the public matrices from the seeds, the LPN sample count is unchanged.

\paragraph{Within-entry row differences.}
A CLPN dictionary entry has rows
\[
  b_i = \langle A_i,r\rangle + x + e_i \pmod q,
\]
where $x$ is the encrypted block inner product.  Subtracting two rows from the same entry cancels $x$:
\[
  b_i-b_j = \langle A_i-A_j,r\rangle + (e_i-e_j) \pmod q.
\]
Thus each dictionary entry yields up to $\binom{m_c}{2}$ message-eliminating q-ary LPN-like rows.  The effective noise rate is estimated by the q-ary piling-up formula with two terms,
\[
  \eta^{(2)} = \frac{q-1}{q}\left(1-\left(1-\frac{q\eta_c}{q-1}\right)^2\right).
\]
The validation code records the raw within-entry surface as
\[
  E\binom{m_c}{2}.
\]

\paragraph{Cross-entry dictionary differences.}
Inside one block, two dictionary entries for coefficient tuples $u$ and $v$ encrypt
\[
  \langle u,s_{\mathrm{block}}\rangle,\qquad \langle v,s_{\mathrm{block}}\rangle.
\]
A row difference between these entries gives
\[
  b_{u,i}-b_{v,j}
  = \langle A_{u,i}-A_{v,j},r\rangle
  + \langle u-v,s_{\mathrm{block}}\rangle
  + e_{u,i}-e_{v,j}.
\]
This is a noisy linear equation in the joint secret $(r,s_{\mathrm{block}})$ of dimension at most $n_c+\max_j\ell_j$.  The raw ordered surface is conservatively screened using
\[
  \sum_j \binom{q^{\ell_j}-1}{2} m_c^2
\]
rows, again with two-term q-ary piling-up noise.  These rows are not all independent, so the repository reports them as a red-flag screening surface rather than a proof of insecurity.

\paragraph{Implication for C3.}
Seeded C3 is still valuable because it avoids materializing $E m_c n_c$ public field elements for dense matrices and keeps only $E m_c$ field elements for the $b$ vectors plus seed metadata.  However, seeded storage does not reduce the within-entry or cross-entry surfaces above.  Therefore the next construction-level problem is not storage but relation control: a publishable C4 variant should avoid a fully closed block dictionary, restrict tuple sets, randomize dictionary availability, or replace the compactor with a code/LPN LHE layer whose public entries do not create so many plaintext-cancelling row differences.

\paragraph{Validation artifact.}
The repository implements this screen in the module
\begin{center}
\path{src/sable/c2_attack_surface.py}.
\end{center}
The generated toy and design reports in the documentation folder show that the current presets remain research-only.  This is expected and useful: the screen identifies the next mathematical bottleneck rather than certifying parameters.
